\section{Methodology}
\label{sec:design}

\subsection{Parallelization Strategy and Decomposition}
We study two ways of using $p$ GPUs.
\emph{Tensor parallelism} (TP) splits one model replica across the GPUs, as described in
\S\ref{sec:background}; it reduces the work and the weight traffic per GPU but adds $2L$
all-reduces per step.
\emph{Data-parallel replicas} (DP) run $p$ independent copies of the model, each serving
$B/p$ requests, with no communication.
Because replicas share nothing, we compute DP throughput as $p$ times the single-GPU
throughput at batch $B/p$, rather than running it separately.
We do not use pipeline parallelism: the model fits on one GPU, and pipeline stages would
add bubbles without reducing per-token latency.

\subsection{Communication and Synchronization}
We measure all-reduce time directly with a microbenchmark that uses the same NCCL library as
\vllm (\texttt{torch.distributed}), for messages from 1\,KB to 256\,MB.
We report bus bandwidth, $\mathit{busbw} = \frac{n}{t}\cdot\frac{2(p-1)}{p}$, the same
definition as \texttt{nccl-tests}.
To expose synchronization and launch overheads we ablate two \vllm features:
(i) the custom all-reduce kernel, replaced by plain NCCL, and
(ii) CUDA graphs, disabled with eager execution.

\subsection{Analytical Performance Model}
We model the time of one step as
\begin{equation}
  T_{\mathrm{step}} = \max\left(T_{\mathrm{comp}},\, T_{\mathrm{mem}}\right) + T_{\mathrm{comm}} .
  \label{eq:step}
\end{equation}
For $N$ tokens in the step ($N=B$ for decode, $N=B\cdot I$ for prefill), with $P$ matrix
parameters, $W$ weight bytes and $K$ bytes of KV cache read,
\begin{equation}
  T_{\mathrm{comp}} = \frac{2PN}{p\,F\,\eta_c}, \qquad
  T_{\mathrm{mem}} = \frac{W + K}{p\,M\,\eta_m},
  \label{eq:compmem}
\end{equation}
where $F$ and $M$ are the peak BF16 throughput and memory bandwidth of one GPU, and
$\eta_c,\eta_m$ are efficiency factors.
The communication term is
\begin{equation}
  T_{\mathrm{comm}} = 2L\left[\alpha_p + \frac{2(p-1)}{p}\cdot\frac{N h s}{\beta_p}\right],
  \label{eq:comm}
\end{equation}
with $s$ bytes per element and per-platform latency $\alpha_p$ and bus bandwidth
$\beta_p$.

\paragraph{Calibration and validation.}
We fit $\alpha_p$ and $\beta_p$ only from the all-reduce microbenchmark, never from \vllm
runs, to avoid circular validation.
We fit $\eta_c$ from TP=1 prefill and $\eta_m$ from TP=1 decode at small batch.
All TP$>$1 points are then \emph{predictions}, and we report their mean absolute
percentage error (MAPE) separately from the calibration set.
The model assumes that communication does not overlap with computation, so it is an upper
bound on step time when overlap exists.

\subsection{Benchmark Harness}
Our harness runs one configuration (model, $p$, engine options) per process.
For each batch size it times a run with one output token (prefill) and a run with $O$
output tokens; the decode step time is
$\mathit{TPOT} = (t_{O} - t_{1})/(O-1)$.
All requests use random token prompts of fixed length and generate exactly $O$ tokens
(end-of-sequence is ignored), and prefix caching is disabled so repetitions do not reuse
cached prompts.

\paragraph{Working under cluster limits.}
The H100 cluster allows jobs of at most 6 minutes, 2 GPUs and 4 CPU cores, with a
per-student budget of about 8 GPU-hours.
The harness therefore (i) keeps the compilation and CUDA-graph caches in persistent
storage so later jobs skip compilation, (ii) appends every measurement to disk
immediately, and (iii) stops itself before a soft time budget, so a killed job still
keeps all completed rows.
The full H100 campaign uses \tbd GPU-minutes.
